Em Mauss e Beuchat, p. 446-447: \enquote{a tenda tem em toda parte o mesmo nome, tupik, e em toda parte também (...) apresenta a mesma forma. (...) Esse tipo normal apresenta naturalmente algumas variações conforme as localidades. (...) Onde a rena é rara, como em Angmagssalik e em toda a Groenlândia oriental, a tenda é feita com peles de foca. Ao mesmo tempo, como ali a madeira não é abundante, a forma da tenda é também um pouco diferente. (...) É curioso observar de que maneira, seja em Igloulik, na baía de Hudson, seja na parte meridional da terra de Baffin, as mesmas causas produsem os mesmos efeitos.}
Em Santos, p. 60: \enquote{cada moradia deveria ser arquitetada com material local, essa é uma primeira grande questão. Todo bioma, todo ambiente, todo lugar nos oferta as condições para viver ali. Se estou nos cocais, posso usar palha das palmeiras para fazer o teto das casas e para forrar as paredes. Se não quiser usar as palhas, posso usar as varas, posso forrar o interior com esteiras, posso também fazer paredes de taipa. Posso fazer minha casa na Caatinga usando apenas o material da Caatinga; posso fazer minha casa no Cerrado usando o material de Cerrado.}